% ══════════════════════════════════════════════════════════════════
% Volume Alpha v2 — 增补稿(P5 交付物)
% 原文 Lingxiao_Xu_Volume_Alpha 不动;本文件为后续版本正文,
% 结果挂接 VolumePrediction/outputs/ 的 CSV/图。
% 编译: xelatex(中文用 Hiragino Sans GB,同 plan PDF 管线)
% 更新: 2026-08-03 — P5 全模型 walk-forward + 窄集 RNN 复赛 + 经济回测 +
%        生产化(numpy 前向服务)成文
% ══════════════════════════════════════════════════════════════════
\documentclass[11pt]{article}
\usepackage{booktabs,graphicx,amsmath,hyperref}
% \usepackage{xeCJK}\setCJKmainfont{Hiragino Sans GB}

\title{Volume Alpha v2: GKMSZ Replication on a PIT R3K Proxy Universe,
       Rolling-Window Model Selection, and a Zero-Dependency
       Production Path for Sequence Models}
\author{Lingxiao Xu}
\date{2026}

\begin{document}\maketitle

\begin{abstract}
We extend the v1 replication in three directions. First, we move from a single
paper split to a twelve-window rolling walk-forward (24m train / 6m test / 6m
step) over a 10.7M-row point-in-time R3K-proxy panel, evaluating ten predictors
under one protocol. Second, we show that the sequence model's advantage reported
in the original Table~1 is \emph{feature-domain conditional}: an LSTM trained on
the narrow technical+size feature set attains a mean out-of-sample
$R^2_\eta$ of 0.295 and beats the best tree model in 11 of 12 windows, while the
same architecture fed all 114 features collapses to 0.139. Third, we translate
the statistical gap into economics under a quadratic-impact accounting protocol
on 2024--2026 out-of-sample predictions, and describe a production path that
serves the sequence model with \emph{no deep-learning dependency at read time}.
\end{abstract}

\section{Data and Universe}          % ← G1/G3;引 P0_data_audit.md 结论
% - PIT R3K 代理构建(§7.3);raw-immutable 缓存+splits 现算
% - 幸存者控制: grouped bar 含退市;FMP 幸存者掩码双轨
% - SIZE 三段拼接+ADV 代理 fallback(ADR 边界)
% - 数据修复链四例(G10 附录):z 截断/横截面 z/字符串股本/重复标签
The production panel (\texttt{prod\_v6}) spans 10,689,240 rows $\times$ 122
columns, 5,620 tickers, with 114 feature columns in five prefix groups
(\texttt{tech\_}, \texttt{fund1\_}, \texttt{fund2\_}, \texttt{cal\_},
\texttt{earn\_}).

\section{Baseline Predictability}    % ← 表: baselines_table.csv;图 fig1
% ma5 93.95 vs paper 93.68 等四项 ±0.6pp

\section{Statistical Models}         % ← 表: table1_{A,B,C};图 fig3

\subsection{Rolling walk-forward over ten predictors}
Table~\ref{tab:wf} reports pooled out-of-sample $R^2_\eta$ across the twelve
windows on the full feature set. Tree and linear models are deterministic
(single fit); the three torch models average five seeds per window.

\begin{table}[h]\centering
\caption{Full-feature-set walk-forward, pooled OOS $R^2_\eta$ (12 windows).}
\label{tab:wf}
\begin{tabular}{lr}
\toprule
Model & $R^2_\eta$ \\
\midrule
LightGBM      & 0.1813 \\
NN2 (legacy)  & 0.1684 \\
NN (paper)    & 0.1572 \\
Forward stepwise & 0.1438 \\
RNN (paper)   & 0.1428 \\
LassoCV       & 0.1394 \\
OLS           & 0.1393 \\
PLS(5)        & 0.1256 \\
PCR(5)        & $-0.0010$ \\
\bottomrule
\end{tabular}
\end{table}

\subsection{The sequence model's advantage is feature-domain conditional}
The v1 Table~1 reported the LSTM leading on the narrow technical+size column
block. Table~\ref{tab:narrow} re-runs the full twelve-window protocol restricted
to that block (14 features) and confirms the effect is not an artifact of the
single paper split: the LSTM gains $+0.117$ mean $R^2_\eta$ over the best
full-set model and wins 11 of 12 windows, whereas the same architecture on all
114 features ranks fifth. Tree models move the opposite way --- LightGBM
\emph{loses} 0.035 when restricted to the narrow block. We read this as a
capacity/regularization interaction: the single-layer LSTM with 32 hidden units
allocates its representational budget to the temporal structure of a small,
homogeneous, well-scaled feature block, and is diluted when that budget must
also absorb 100 heterogeneous cross-sectional columns.

\begin{table}[h]\centering
\caption{Narrow (tech+fund1, 14 features) vs full (114 features).}
\label{tab:narrow}
\begin{tabular}{lrr}
\toprule
Model & narrow & full \\
\midrule
RNN  & \textbf{0.2947} & 0.1388 \\
LightGBM & 0.1428 & 0.1773 \\
\bottomrule
\end{tabular}
\end{table}

\subsection{A harness pitfall specific to sequence models}
Evaluating a many-to-one sequence model window-by-window introduces a subtle
penalty absent for cross-sectional models: the first $L-1$ rows of each ticker
inside a test window have their lookback truncated at the train/test boundary
and are zero-padded, even though that history is strictly past data already
available to the model. At $L=10$ and 126-day test windows this affects 7.1\% of
evaluation rows. Supplying the pre-boundary context (strictly earlier than the
first test date, hence no look-ahead) is required for a like-for-like comparison
with cross-sectional models; all sequence results reported here include it.

\section{Economic Framework}         % ← 表: table2_mel/table3_transfer;图 fig4
% 迁移签名(d_mse>0 与 MEL 增益并存);
% 从零直训失效 → G7_divergence_report 摘编为 v2 的一节(新贡献)

\section{Trading Experiments}        % ← 图 fig5/fig7;表 table3_trading

\subsection{Does the statistical gap pay?}
We evaluate 2024-01 through 2026-07 (644 trading days) under the accounting
protocol of the v1 economic section: exogenous oracle trade demand, quadratic
impact $\lambda\,(\text{traded})^2/V$, and the closed-form participation rule
$s(\hat v;\mu)$, so that tiers differ \emph{only} through forecast quality. All
model forecasts are genuine out-of-sample walk-forward predictions.
Table~\ref{tab:econ} reports each tier at its Sharpe-optimal $\mu$.

\begin{table}[h]\centering
\caption{Net performance, 2024--2026, at Sharpe-optimal $\mu=10^{-9}$.}
\label{tab:econ}
\begin{tabular}{lrrr}
\toprule
AUM & Tier & Net ann.\ ret.\ (\%) & Net Sharpe \\
\midrule
\$1e8  & RNN (narrow) & \textbf{62.19} & \textbf{7.00} \\
       & LightGBM (full) & 60.23 & 6.19 \\
       & MA5 & 58.58 & 6.25 \\
\$1e9  & RNN (narrow) & \textbf{57.58} & \textbf{6.48} \\
       & LightGBM (full) & 55.73 & 5.73 \\
       & MA5 & 53.98 & 5.76 \\
\$1e10 & RNN (narrow) & \textbf{11.44} & \textbf{1.28} \\
       & LightGBM (full) & 10.76 & 1.10 \\
       & MA5 & 7.99 & 0.85 \\
\bottomrule
\end{tabular}
\end{table}

At \$1B the sequence model adds 1.85pp of net annual return and 0.76 of Sharpe
over the incumbent tree model; at \$10B, where cost drag exceeds 50\%, it adds
3.4pp over the naive MA5 forecast. We note that at \$1e8--\$1e9 all forecast
tiers post higher net Sharpe than the oracle tier: knowing $v$ exactly leads the
participation rule to schedule more aggressively, raising both turnover and the
variance of realized cost. This is a property of the accounting protocol rather
than of the forecasts, and does not affect comparisons among the three
predictors, which face identical demand and cost parameters.

\section{Legacy Model Zoo on the New Panel}  % ← legacy_scores.csv
% 12 模型并排+旧成绩对照;定性排序验收

\section{Productionization}          % ← service.py/profiles/registry

\subsection{Serving a sequence model without a deep-learning runtime}
The read path of the production service is forbidden from importing torch: the
daily job must start in seconds, and a training-framework dependency in the
serving environment is an availability liability. Since the deployed network is
small --- a single LSTM layer with 32 hidden units over 14 inputs plus three
dense layers, $\approx 6.7$k parameters --- we export the trained weights and
reimplement the forward pass in NumPy (roughly forty lines: the standard
$i,f,g,o$ gate decomposition plus three dense layers). Agreement with the torch
forward pass is $1.5\times10^{-8}$ in maximum absolute deviation, and a test
blocks \texttt{import torch} in a subprocess to verify the serving path runs
without it.

\subsection{Stateful serving}
Unlike the cross-sectional models, a many-to-one sequence model requires the
last $L-1$ feature rows per ticker at serve time. The frozen artifact therefore
carries a \texttt{seq\_tail} tensor (\#tickers $\times$ 9 $\times$ 14) with a
date stamp; the daily job appends the current row, drops the oldest, and
rewrites atomically. Three disciplines are enforced by tests: repeated
same-day serving is idempotent (stamp comparison), a date gap raises rather than
silently predicting from a misaligned window, and predictions are averaged over
the frozen seeds rather than trusting a single training run.

\subsection{Other production notes}
% 六 utility profile;mu_calibration(aiss 2.82e-3 实测);
% walkforward 分层(行业 spread 13pp);日更影子验收(P3/P4 追加)
Stratified walk-forward reporting (industry and size decile) accompanies every
run; industry $R^2$ spread reaches 12pp for the incumbent model, which is the
basis for the per-strategy utility profiles.

\appendix
\section{Divergence Reports}         % G7 全文
\section{58-Step Checklist}          % appendixA_checklist.md 摘要
\end{document}
